% -*- program: xelatex -*-
%\documentclass[12pt, compress]{beamer} 
\documentclass[11pt, compress]{beamer}       
%\documentclass[notes=only]{beamer}

\usepackage{fontspec}
	
\setsansfont{PT Sans}
\setmonofont{Liberation Mono}

\usepackage{beamerthemesplit} 
\usefonttheme[onlymath]{serif}
%\DeclareFontEncoding{T1}
\fontencoding{T1}

%\usepackage[T2A]{fontenc} 
\usepackage[utf8x]{inputenc}
\usepackage[russian]{babel}
\usepackage{listings}

\usepackage{amsmath,amssymb,amsthm}
\hypersetup{unicode=true}
\usepackage{graphics,graphicx}
\usepackage{verbatim}
\usepackage{fancyvrb}
\usepackage{booktabs}
%\usepackage{media9}
\usepackage[super]{natbib}
\usepackage{color}
\definecolor{mygreen}{RGB}{28,172,0} % color values Red, Green, Blue
\definecolor{mylilas}{RGB}{170,55,241}
\definecolor{codecolor}{RGB}{0,120,0}
\definecolor{light-gray}{gray}{0.95}
\definecolor{light-green}{rgb}{0.93, 1, 0.8}
\definecolor{light-yellow}{rgb}{1,1,0.85}
\definecolor{dark-blue}{rgb}{0,0,0.6}
\definecolor{dark-red}{rgb}{0.7,0,0.0}
\definecolor{dark-green}{RGB}{10,150,0}

%\usepackage{pifont}% http://ctan.org/pkg/pifont
\newcommand{\cmark}{\textcolor{green}{\ding{51}}}
\newcommand{\xmark}{\textcolor{red}{\ding{55}}}

\setbeamercolor{frametitle}{fg=dark-red,bg=gray!10}
\setbeamerfont{title}{series=\bfseries,parent=structure}
\setbeamerfont{frametitle}{series=\bfseries,parent=structure}

\setbeamertemplate{blocks}[rounded][shadow=false]
\setbeamertemplate{navigation symbols}{}

\newcommand{\code}[1]{\textcolor{dark-green}{\texttt{#1}}}
\newcommand{\cleartitlepage}[1]{\begin{frame}[plain]\begin{center}\textsc{#1}\end{center}\end{frame}}
\renewcommand{\emph}[1]{\textcolor{dark-blue}{#1}}
\newcommand{\emphb}[1]{\textcolor{dark-blue}{\textbf{#1}}}
\newcommand{\lstcomment}[1]{\textcolor{dark-green}{\# \texttt{#1}}}

\usepackage{dirtree}

\usetheme{CambridgeUS}
\usecolortheme{lily}

\lstset{basicstyle=\ttfamily,
    language=python,    
    keepspaces=true,
    extendedchars=\true,
    basicstyle={\small},
    breaklines=true,
    morekeywords={matlab2tikz},
    keywordstyle=\color{blue},
    morekeywords=[2]{1}, keywordstyle=[2]{\color{blue}},
    identifierstyle={ \bf \color{black} },
    stringstyle=\color{mylilas},
    commentstyle=\color{mygreen},%
    showstringspaces=false,
    numbers=left,%
    numberstyle={\scriptsize \color{gray}},
    numbersep=7pt, 
    emph=[1]{case,switch,otherwise,nonlocal,as,yield, with},emphstyle=[1]\color{blue}, 
    frame=single,    
    rulecolor=\color{red},    
    backgroundcolor=\color{light-gray},
    xleftmargin=0.3cm,
    frame=l,framesep=4pt,framerule=0.5pt,
    escapechar=|
    %lineskip=-1.0pt,
    %emph=[2]{word1,word2}, emphstyle=[2]{style},    
}

\usepackage{tikz}
\usepackage{pgfplots}

\usetikzlibrary{positioning} 
\tikzset{cbutton/.style={rectangle,minimum width=10mm,thick,draw=red!50!black!50,
top color=white,bottom color=red!50!black!30}}
\tikzset{mbutton/.style={rectangle,minimum width=10mm,thick,draw=black!20,
top color=white,bottom color=black!30}}

\usetikzlibrary{arrows,shapes}
\tikzstyle{every picture}+=[remember picture]


\makeatother
\setbeamertemplate{footline}
{
  \leavevmode%
  \hbox{%  
  \begin{beamercolorbox}[wd=.3\paperwidth,ht=2.25ex,dp=1ex,center]{author in head/foot}%
    \usebeamerfont{author in head/foot}\insertshortauthor
  \end{beamercolorbox}%  
  \begin{beamercolorbox}[wd=.6\paperwidth,ht=2.25ex,dp=1ex,center]{title in head/foot}%
    \usebeamerfont{title in head/foot}\insertshorttitle
  \end{beamercolorbox}%
  \begin{beamercolorbox}[wd=.1\paperwidth,ht=2.25ex,dp=1ex,center]{date in head/foot}%
    \insertframenumber{} / \inserttotalframenumber\hspace*{1ex}
  \end{beamercolorbox}}%
  \vskip0pt%
}

\setbeamertemplate{headline}{}
\setbeamertemplate{footline}{}
 
\AtBeginSection[]{
  \begin{frame}[plain]
  %\vfill  
  \begin{tikzpicture}
  \useasboundingbox (0,0) rectangle(\the\paperwidth,\the\paperheight);  
  \fill[color=dark-red]   (-1cm, 3.9cm) rectangle(\the\paperwidth, 5.9cm);
  %\usebeamerfont{title}\insertsectionhead\par%
  %\node[text width=\the\paperwidth,align=center] at (current page.center) {\color{ExecusharesWhite}\Large\textbf{\insertsectionhead}};  
  \node[text width=\the\paperwidth,align=center] at (6cm,4.9cm) {\color{white}\Large\textbf{\insertsectionhead}};  
  \end{tikzpicture}
  \end{frame}
}

% ============================================================================

\title[Python: ввод и вывод]{Ввод и вывод}
\subtitle{Технологии и языки программирования}
\author[Самарский университет]{Юдинцев В. В.}
\institute{Кафедра теоретической механики\\Самарский университет}
\date{\today}


\begin{document}

{
\setbeamercolor{background canvas}{bg=light-yellow} 
\usebackgroundtemplate[background]{}
%\begin{frame}[plain]
%\begin{figure}[h]
%  \centering
%  \includegraphics[width=0.3\textwidth]{../common/Python_logo_and_wordmark.pdf}
%\end{figure}
%\maketitle
%\end{frame}
}

% \frame{\frametitle{Содержание}\tableofcontents}

%\section{Форматированный вывод}

% axis style, ticks, etc
\pgfplotsset{every axis/.append style={
                    axis x line=middle,    % put the x axis in the middle
                    axis y line=middle,    % put the y axis in the                     
                    xlabel={$x$},          % default put x on x-axis
                    ylabel={$y$},          % default put y on y-axis
                    label style={font=\small},
                    tick label style={font=\small}  
                    }}

\begin{frame}[c, fragile]
\frametitle{Действие постоянной силы}
\begin{columns}
\column{0.45\linewidth}
\begin{tikzpicture} 
\usetikzlibrary{decorations.pathmorphing,patterns}
\tikzstyle{spring}=[thick,decorate,decoration={zigzag,pre length=0.1cm,post
  length=0.1cm,segment length=6}]
\draw[spring] (0.5,0) -- (4,0);  
\draw[thick,red,->] (4,0) -- (5,0) node[anchor=south east] {$\boldsymbol P$};  
\draw[thick,red,->] (4,0) -- (2.5,0) node[anchor=south] {$\boldsymbol F = c x$};  
\node[circle,fill=blue,inner sep=1.5mm, text=white] (a) at (4,0) {m};
\fill [pattern = north east lines] (0,-1) rectangle (0.5,1);
\draw [->] (0.5,-0.6) -- (4,-0.6) node[anchor = north east] {$x$};
\draw [-] (4,-0.7) -- (4,-0.4);
\end{tikzpicture}
\begin{tikzpicture} 
\begin{axis}[xlabel=,ylabel=, height=6cm, width=5cm,
		ymin=0,ymax=150,
		ytick = {0,100,150},
		yticklabels = {0, $P$, },		
		xmin=0,xmax=2,
		xtick = {0,1,2},
      	xticklabels = {0, $\tau$, t}] 
\addplot[color=red,mark=o] coordinates { (0,100) (1,100) (1,0)}; 
\end{axis} 
\end{tikzpicture}
\column{0.55\linewidth}
Зависимость упругой силы $F$ от продолжительности действия внешней силы $P$: 
$$
	F = P \sqrt{ 2-2 \cos \omega \tau }
$$
При 
$$
\tau = \frac{\pi}{\omega} = \pi \sqrt{\frac{m}{c}} ,
$$
т. е. при продолжительности действия силы $P$ равной половине периода собcтвенных колебаний системы груз-пружина ($T=2 \pi / \omega$), максимальное усилие в пружине в два раза больше силы P
$$
	F(\tau = \pi \sqrt{m/c} ) = 2 P
$$
\end{columns}
 
\end{frame}

\begin{frame}
\frametitle{Решение}
Уравнение движения системы при действии силы $P$
\begin{equation}
	m \ddot{x} = P - c x, 	
\end{equation}
Начальные условия
\begin{equation}
	x_0 = 0, \quad \dot{x}_0 = 0
\end{equation}
Решение
\begin{equation}
	x = c_1 \sin \omega t + c_2 \cos \omega t + P/c, \quad \omega = \sqrt{c/m}
\end{equation}
Для начальных условий $c_1 = 0$, $c_2 = - P/c$ и 
\begin{equation}
	x = \frac{P}{c}\left( 1-\cos \omega t \right), \quad \dot x = \frac{P \omega }{c} \sin \omega t 
\end{equation}
\end{frame}

\begin{frame}[t]
Полная энергия системы в момент окончания действия силы $P$
\begin{equation}
	E_\tau = \frac{c x_\tau^2}{2} + \frac{m \dot{x}_\tau^2}{2}
\end{equation}
Полная энергия системы в момент максимальной деформации пружины
\begin{equation}
	E = \frac{c x_{max}^2}{2} 
\end{equation}
Учитывая, что 
$$
	\frac{c x_\tau^2}{2} + \frac{m \dot{x}_\tau^2}{2} = \frac{c x_{max}^2}{2} 
$$
Максимальная деформация
\begin{equation}
	x_{max}^2 = \frac{P^2}{c^2}\left( 1-\cos \omega \tau \right)^2+\frac{P^2} {c^2} \sin^2 \omega \tau  
\end{equation}
Максимальная сила
\begin{equation}
	F = c x_{max} = P \sqrt{2 (1-\cos \omega \tau)}
\end{equation}

	
\end{frame}


\end{document}